\begin{answers}{ch:rl-basics}{強化学習の基礎：MDP と価値関数}
\item $Q^\pi(6,\text{上})=0.0580$ です。
状態 6 で「上」を選ぶと、確率 0.8 で状態 2（報酬 0）、0.1 で壁にぶつかって状態 6 のまま（報酬 0）、0.1 で穴 7（報酬 $-1$）に移ります。
穴は終端状態なので、その先の価値は 0 です。
\eqref{eq:rlb-q-from-v}に代入すると、次のようになります。
\begin{align*}
Q^\pi(6,\text{上})&=0.8\,(0+0.9\times0.2616)+0.1\,\bigl(0+0.9\times(-0.3371)\bigr)+0.1\,(-1+0)\\
&=0.1884-0.0303-0.1=0.0580.
\end{align*}
\item 状態も次の状態も同じ 1 つなので、\eqref{eq:rlb-bellman-v}の和はすべて 1 項になり、$V=1+\gamma V$ です。
どの行動の確率を足しても 1、遷移確率も 1 だからです。
$V$ を左辺に集めると $(1-\gamma)V=1$、$0\le\gamma<1$ なので $V=1/(1-\gamma)$ です。
収益 $1+\gamma+\gamma^2+\cdots$ の等比級数の和と一致します（$\gamma=0.9$ なら 10）。
\item $\gamma=0.5$ では、状態 9 が「左」から「右」に、状態 11 が「左」から「下」に変わります（実測。$V^\ast(8)=0.0331$）。
割引が強いと、遠回りの上の道で着くゴールの価値が小さくなり、穴の横を通っても短い道（9→10→6→2）のほうが得になります。
状態 11 では、左へ進むと 0.1 の確率で穴に滑ります。
その先の価値が小さいので危険に見合わず、壁に向かって「下」を選び、滑ったときだけ左へ出る安全な行動が最適になります。
1 歩ごとに $-0.04$ にすると、状態 9 だけが「右」に変わります（実測。$V^\ast(8)=0.3739$）。
1 歩ごとに罰があるので、多少の危険をおかしても早く着くほうが得になるからです。
状態 11 は「左」のままです。
\item SARSA の更新式は $Q(s_t,a_t)\leftarrow Q(s_t,a_t)+\alpha\,[\,r_t+\gamma Q(s_{t+1},a_{t+1})-Q(s_t,a_t)\,]$ です。
コードでは、$a_{t+1}$ を先に選んでおき、次の 1 歩でその行動を実際に使うように書き換えます。
結果は、SARSA の $Q$ が全体に小さくなり、貪欲方策はほぼ同じでした（実測。$\alpha=0.1$、$\varepsilon=0.1$、2,000 エピソード）。
例えば $Q(2,\text{右})$ は Q 学習の 0.990 に対して 0.867、$\max_a Q(8,a)$ は 20 回の平均で 0.531 に対して 0.497 です。
SARSA は「10\% はでたらめに動く」方策そのものの価値を学ぶので、探索中に穴へ落ちる損が $Q$ に入ります。
貪欲方策は、どちらも 9 状態のうち 8 状態で最適方策と一致しました（違うのは、Q 学習が状態 11、SARSA が状態 10 で、どちらも訪れにくい右下です）。
探索が穴に近い道を危険にするので、SARSA は穴の近くで安全側の行動（状態 10 で「左」）を選ぶことがあります（20 回中 13 回）。
\item 状態はプロンプト「12+35は？」とここまでに生成したトークン列、行動は次の 1 トークン（語彙の 4,096 通りのどれか）です。
遷移は選んだトークンを末尾に足すだけなので、確率 1 で決まります。
報酬は、答え終わったときに正しければ 1、間違いなら 0 を 1 回だけ与え、途中の報酬は 0 とするのが自然です。
エピソードは、終わりの記号 \code{<|end|>} を書いたとき（または長さの上限に達したとき）に終わります。
方策は、LLM が出す次のトークンの確率分布 $\pitheta(a\mid s)$ そのものです。
自分で答えてから、\tabref{tab:rlb-maze-llm}と見比べてください。
表の「集めるデータ」「学習するもの」「状態の再訪」の行も、格子世界との違いとして確かめておくと、\chapref{ch:rl-llm}が読みやすくなります。
\end{answers}
